\begin{courseunit}{I.1}
\chapter{Unit 1: The Grammar Bridge}

\begin{worksheet}{I.1}{A}{Unit 1: The Grammar Bridge}{Variant A}{Clause roles, dictionary forms, and first reading habits}{Use independently after the unit; write every requested label, not only a translation.}
\Exercise{Copy \Latin{Dominus custodit animam}. Underline the subject once, the finite verb twice, and box the direct object. Then give the neutral English order.}{Subject: \Latin{Dominus}; finite verb: \Latin{custodit}; direct object: \Latin{animam}. \English{The Lord guards the soul.}}
\Exercise{Label each entry noun, verb, or adjective and state what its listed parts tell you: \Latin{Deus, Dei m.; gratia, gratiæ f.; laudo, laudare, laudavi, laudatum; sanctus, -a, -um}.}{\Latin{Deus} and \Latin{gratia} are nouns; nominative, genitive, and gender are supplied. \Latin{Laudo} is a verb; its four principal parts are supplied. \Latin{Sanctus} is an adjective; its masculine, feminine, and neuter nominative singular forms are supplied.}
\Exercise{Rewrite \Latin{Deum laudamus} in subject--verb--object order. Explain why the meaning does not change.}{\Latin{Nos laudamus Deum}. With the normally omitted subject pronoun, \Latin{Laudamus Deum} preserves the verb--object sequence but no longer displays all three requested positions. Verb ending \Latin{-mus} identifies \English{we}; \Latin{Deum} remains the object regardless of position.}
\Exercise{In \English{The eternal Word guards our souls}, name the subject, finite verb, object, and every modifier.}{Subject: \English{Word}; finite verb: \English{guards}; object: \English{souls}; \English{the} and \English{eternal} modify \English{Word}; \English{our} modifies \English{souls}.}
\Exercise{Translate \Latin{Deo gratias}. State what may be supplied to make a
full Latin clause.}{\English{Thanks be to God.} \Latin{Deo} is the recipient;
\Latin{gratias} means \English{thanks}. A close Latin expansion may supply
\Latin{[agimus]} or \Latin{[agamus]}. The word \English{be} belongs to the
conventional English rendering and is not an omitted Latin copula.}
\PassageExercise{In principio erat Verbum.}{Mark the prepositional phrase, finite verb, and subject; then translate.}{Prepositional phrase: \Latin{in principio}; finite verb: \Latin{erat}; subject: \Latin{Verbum}. \English{In the beginning was the Word.}}
\end{worksheet}

\begin{worksheet}{I.2}{B}{Unit 1: The Grammar Bridge}{Variant B}{Clause roles, orthographic conventions, and first reading habits}{This is interchangeable with Variant A; complete it without consulting that sheet.}
\Exercise{Analyze \English{Holy light fills the church}: identify subject, finite verb, object, and modifiers.}{Subject: \English{light}; finite verb: \English{fills}; object: \English{church}; \English{holy} modifies \English{light}; \English{the} modifies \English{church}.}
\Exercise{For \Latin{pax, pacis f.; audio, audire, audivi, auditum; omnis, omne}, identify part of speech and the minimum information needed for lookup.}{\Latin{Pax} is a noun: nominative, genitive, gender. \Latin{Audio} is a verb: four principal parts. \Latin{Omnis} is a two-termination adjective: masculine/feminine and neuter nominative singular; its genitive is found under the entry.}
\Exercise{Normalize these variant forms to this course's 1962 house spelling: \Latin{JESVS, EVANGELIVM, praeceptum}. State the three graphic conventions involved.}{\Latin{Iesus, Evangelium, præceptum}. Another edition may use \Latin{J} for consonantal \Latin{i}, capital \Latin{V} for both \Latin{u/v}, or expand ligature \Latin{æ} as \Latin{ae}. Capitalization may also differ.}
\Exercise{Copy \Latin{Gloria Patri et Filio}. Separate the coordinated terms and state the repeated relationship expressed by their endings.}{The coordinated terms are \Latin{Patri} and \Latin{Filio}. Both name recipients: \English{Glory to the Father and to the Son.}}
\Exercise{Translate \Latin{Dominus vobiscum}. Put brackets around any interpretive expansion.}{\English{The Lord [be] with you}; the conventional optative construal may be shown as \Latin{Dominus [sit] vobiscum}, but \Latin{sit} is neither printed nor uniquely required by morphology.}
\PassageExercise{Dominus vobiscum. Et cum spiritu tuo.}{Identify which response words form one prepositional phrase and give a natural translation of the exchange.}{\Latin{cum spiritu tuo} is one phrase. \English{The Lord be with you. And with your spirit.} The second response is elliptical rather than a complete printed clause.}
\end{worksheet}

\begin{worksheet}{I.3}{C}{Unit 1: The Grammar Bridge}{Variant C}{Inflection, English grammar, and disciplined observation}{Work from visible endings and context; do not guess from familiar translations.}
\Exercise{In \English{The priest gives the people peace}, identify subject, finite verb, indirect object, and direct object.}{Subject: \English{priest}; finite verb: \English{gives}; indirect object/recipient: \English{people}; direct object: \English{peace}.}
\Exercise{Copy \Latin{Animam custodit Dominus}. Rearrange it twice while preserving meaning and explain what English must use instead of endings.}{Possible forms: \Latin{Dominus custodit animam}; \Latin{Custodit Dominus animam}. English normally relies on fixed word order: \English{The Lord guards the soul}; changing English noun order would change roles.}
\Exercise{Expand the analysis checklist for \Latin{laudamus}: write the questions to ask before translating it.}{Model: What is the dictionary verb and conjugation? Is the form finite? Which person and number does the ending mark? Which tense, voice, and mood are present? What is its subject, stated or understood? Does it govern an object or complement? Only then select an English rendering.}
\Exercise{Explain why the entry \Latin{rex, regis m.} cannot safely be learned as \Latin{rex} alone.}{\Latin{Regis} reveals the stem \Latin{reg-} and declensional pattern; \Latin{m.} gives gender. The nominative \Latin{rex} alone does not predict forms such as \Latin{regem} or \Latin{regibus}.}
\Exercise{Translate \Latin{Agnus Dei} closely and identify whether it is a full clause.}{\English{Lamb of God.} It is a noun phrase/acclamation, not a full clause: no finite verb is printed.}
\PassageExercise{Et Verbum caro factum est.}{Identify the conjunction, the two noun-like forms, and the complete finite verb group. Give a provisional translation.}{Conjunction: \Latin{et}. Noun-like nominatives: \Latin{Verbum} and \Latin{caro}. Finite group: \Latin{factum est}. \English{And the Word was made flesh} or \English{And the Word became flesh.}}
\end{worksheet}

\begin{worksheet}{I.4}{D}{Unit 1: The Grammar Bridge}{Variant D}{Dictionary literacy, word order, and close translation}{Use different words to revisit the same abilities.}
\Exercise{Analyze \English{Our Father grants eternal life to the faithful}: identify subject, finite verb, direct object, recipient, and modifiers.}{Subject: \English{Father}; finite verb: \English{grants}; direct object: \English{life}; recipient: \English{the faithful}; \English{our} modifies \English{Father}; \English{eternal} modifies \English{life}.}
\Exercise{Classify the listed information in \Latin{corpus, corporis n.; credo, credere, credidi, creditum; æternus, -a, -um}.}{\Latin{Corpus} is a neuter noun with nominative and genitive; \Latin{credo} is a verb with four principal parts; \Latin{æternus} is a three-termination adjective giving masculine, feminine, and neuter nominative singular.}
\Exercise{For printed \Latin{GLORIA IN EXCELSIS DEO}, write the course form and state why capitalization does not supply grammatical roles.}{\Latin{Gloria in excelsis Deo}. Capitals are a typographic choice; the endings and syntax, not letter case, show relationships.}
\Exercise{Copy \Latin{Te laudamus}. Supply its unprinted subject, name the object, and translate.}{The verb ending supplies subject \Latin{nos}, \English{we}; \Latin{te} is the object. \English{We praise you.}}
\Exercise{Diagnose this reading error: ``The first noun in \Latin{Deum laudat Maria} must be the subject.'' Give the correct analysis.}{Latin need not put the subject first. \Latin{Deum} is the object, \Latin{laudat} is the finite verb, and \Latin{Maria} is the subject: \English{Mary praises God.}}
\PassageExercise{Deo gratias. Amen.}{This course composite places two independent liturgical responses together; separate their syntax and give each sense.}{\Latin{Deo gratias} is an elliptical expression, \English{Thanks be to God}; \Latin{Amen} is an indeclinable acclamation, \English{Amen} or \English{truly/so be it}. They are not presented as one continuous received source passage.}
\end{worksheet}

\end{courseunit}

\begin{courseunit}{I.2}
\chapter{Unit 2: Case and Agreement}

\begin{worksheet}{I.5}{A}{Unit 2: Case and Agreement}{Variant A}{Case functions, agreement, and ambiguity resolved by syntax}{Give case, number, function, and translation where requested.}
\Exercise{Match each case to its core function: nominative, genitive, dative, accusative, ablative, vocative; subject, possession, recipient, direct object, circumstance/prepositional object, direct address.}{Nominative--subject; genitive--possession or description; dative--recipient; accusative--direct object; ablative--circumstance or object of an ablative-governing preposition; vocative--direct address.}
\Exercise{Parse \Latin{Patri} and \Latin{Filio} in \Latin{Gloria Patri, et Filio}: state all four coordinates and function.}{Both are masculine singular dative and recipients of the wished glory: \English{to the Father and to the Son}.}
\Exercise{In \Latin{anima bona}, change both words together to accusative singular and nominative plural.}{Accusative singular: \Latin{animam bonam}. Nominative plural: \Latin{animæ bonæ}. The adjective changes because it must match gender, number, and case.}
\Exercise{Transform \Latin{Dominus custodit animam}, \English{the Lord guards the soul}, into \English{the soul guards the Lord}.}{\Latin{Anima custodit Dominum}. The noun endings, not merely position, exchange subject and object.}
\Exercise{Diagnose \Latin{Deus laudat gloria} if the intended meaning is \English{God praises glory}; correct and explain.}{\Latin{Deus laudat gloriam}. \Latin{Gloriam} must be accusative singular because it is the direct object; \Latin{gloria} would normally be nominative or ablative singular.}
\PassageExercise{Pax hominibus bonæ voluntatis.}{Parse the three nouns and translate, marking any interpretive expansion.}{\Latin{Pax}: feminine singular nominative in a verbless acclamation or wish. \Latin{hominibus}: masculine plural dative, recipient. \Latin{voluntatis}: feminine singular genitive, dependent on \Latin{hominibus}; \Latin{bonæ} agrees with it. Conventional English supplies \English{be}; analytic Latin may mark interpretive \Latin{[sit]}. \English{Peace [be] to people of good will.}}
\end{worksheet}

\begin{worksheet}{I.6}{B}{Unit 2: Case and Agreement}{Variant B}{Case functions, agreement, and direct address}{This variant is equivalent in scope to Variant A.}
\Exercise{For each bold idea in \English{O Lord, give peace to the people of the Church}, name the expected Latin case: Lord, peace, people, Church.}{\English{Lord}: vocative; \English{peace}: accusative direct object; \English{people}: dative recipient; \English{Church}: genitive possession or association.}
\Exercise{Parse every declined word in \Latin{Domine, non sum dignus}.}{\Latin{Domine}: masculine singular vocative, direct address. Understood subject of \Latin{sum}: first-person singular. \Latin{dignus}: masculine singular nominative predicate adjective agreeing with the understood speaker.}
\Exercise{Make adjective \Latin{æternus, -a, -um} agree with (a) \Latin{vitam}, (b) \Latin{lumine}, and (c) \Latin{sæcula}.}{(a) \Latin{vitam æternam}, feminine singular accusative; (b) \Latin{lumine æterno}, neuter singular ablative; (c) \Latin{sæcula æterna}, neuter plural nominative or accusative.}
\Exercise{Explain both possible roles of isolated \Latin{Deo}, then resolve it in \Latin{Deo gratias}.}{\Latin{Deo} can be masculine singular dative or ablative. In \Latin{Deo gratias} it is dative recipient, \English{thanks to God}; no ablative-governing word or circumstance is present.}
\Exercise{Diagnose \Latin{cum Domino æternum}; correct the adjective and state the governing rule.}{\Latin{cum Domino æterno}. \Latin{Cum} governs ablative; \Latin{Domino} is masculine singular ablative, so its adjective must be \Latin{æterno}.}
\PassageExercise{Gloria Patri, et Filio, et Spiritui Sancto.}{Identify the implied element, parse the coordinated recipients, and translate.}{A wish such as \Latin{sit}, \English{be}, is implied. \Latin{Patri, Filio}, and \Latin{Spiritui Sancto} are masculine singular datives; \Latin{Sancto} agrees with \Latin{Spiritui}. \English{Glory be to the Father, and to the Son, and to the Holy Spirit.}}
\end{worksheet}

\begin{worksheet}{I.7}{C}{Unit 2: Case and Agreement}{Variant C}{Case selection from clause function and government}{Show the question that led to each case choice.}
\Exercise{Assign a case to each role in \English{The priest gives the Lord's book to Mary}: priest, Lord's, book, Mary.}{\English{priest}: nominative subject; \English{Lord's}: genitive possessor; \English{book}: accusative object; \English{Mary}: dative recipient.}
\Exercise{In \Latin{Agnus Dei, qui tollis peccata mundi}, identify the case and function of both genitives and the direct object.}{\Latin{Dei}: masculine singular genitive modifying \Latin{Agnus}. \Latin{mundi}: masculine singular genitive modifying \Latin{peccata}. \Latin{peccata}: neuter plural accusative direct object of \Latin{tollis}.}
\Exercise{Change \Latin{lux perpetua} from nominative singular to genitive singular, dative singular, and accusative singular.}{Genitive: \Latin{lucis perpetuæ}; dative: \Latin{luci perpetuæ}; accusative: \Latin{lucem perpetuam}.}
\Exercise{Why can \Latin{sanctis} not be assigned one case without context? Give its possible coordinates as an adjective form.}{It may be dative or ablative plural in any gender. A governing preposition, recipient role, or noun agreement must resolve it.}
\Exercise{Correct \Latin{in terra sanctam} when it means \English{on holy earth}; explain both corrections or confirmations.}{\Latin{in terra sancta}. Here \Latin{in} denotes place where and governs the ablative; \Latin{terra} is already feminine singular ablative, and adjective \Latin{sancta} must agree.}
\PassageExercise{Agnus Dei, qui tollis peccata mundi, dona nobis pacem.}{Parse \Latin{Agnus, Dei, peccata, mundi, nobis}, and \Latin{pacem}; then translate.}{\Latin{Agnus}: nominative singular used for direct address; \Latin{Dei}: genitive singular; \Latin{peccata}: accusative plural; \Latin{mundi}: genitive singular; \Latin{nobis}: dative plural recipient; \Latin{pacem}: accusative singular object. \English{Lamb of God, who take away the sins of the world, grant us peace.}}
\end{worksheet}

\begin{worksheet}{I.8}{D}{Unit 2: Case and Agreement}{Variant D}{Case contrasts, coordinated phrases, and correction}{Do not use word position as proof of case.}
\Exercise{Give the Latin case expected for each italicized relation: glory \emph{to God}; Body \emph{of Christ}; \emph{with the spirit}; hear \emph{the word}; \emph{O Father}.}{Dative; genitive; ablative after \Latin{cum}; accusative direct object; vocative.}
\Exercise{Parse \Latin{animam meam} and \Latin{Corpus Domini} by coordinates and relation.}{\Latin{animam meam}: feminine singular accusative noun and agreeing possessive, an object phrase. \Latin{Corpus}: neuter singular nominative or accusative as context decides; \Latin{Domini}: masculine singular genitive modifying it, \English{Body of the Lord}.}
\Exercise{Transform \Latin{Deus bonus animam custodit} so that plural gods guard plural souls.}{\Latin{Dei boni animas custodiunt}. Nominative subject, adjective, accusative object, and verb number all change.}
\Exercise{Explain why \Latin{Patris} and \Latin{Patri} are not interchangeable in \English{the name of the Father} and \English{glory to the Father}.}{\Latin{Patris} is genitive singular for possession: \Latin{nomen Patris}. \Latin{Patri} is dative singular for recipient: \Latin{gloria Patri}.}
\Exercise{Diagnose and correct \Latin{dona nos pacem}, intended \English{grant us peace}.}{\Latin{Dona nobis pacem}. Recipient \English{us} requires dative \Latin{nobis}; direct object \Latin{pacem} remains accusative.}
\PassageExercise{Corpus Domini nostri Iesu Christi custodiat animam meam.}{Draw the genitive and agreement links, identify subject and object, and translate.}{Subject: neuter singular nominative \Latin{Corpus}. Genitive name chain \Latin{Domini nostri Iesu Christi} depends on it; \Latin{nostri} agrees with \Latin{Domini}. Object: feminine singular accusative \Latin{animam meam}. \English{May the Body of our Lord Jesus Christ preserve my soul.}}
\end{worksheet}

\end{courseunit}

\begin{courseunit}{I.3}
\chapter{Unit 3: First and Second Declensions}

\begin{worksheet}{I.9}{A}{Unit 3: First and Second Declensions}{Variant A}{Paradigms, neuters, vocatives, and liturgical noun phrases}{Write full forms, not ending fragments.}
\Exercise{Give these forms of \Latin{missa, missæ f.}: genitive singular, dative plural, accusative plural, ablative singular.}{\Latin{missæ, missis, missas, missa}.}
\Exercise{Give these forms of \Latin{dominus, domini m.}: vocative singular, accusative singular, nominative plural, genitive plural.}{\Latin{domine, dominum, domini, dominorum}.}
\Exercise{Apply the neuter rule to \Latin{verbum, verbi n.}: give nominative/accusative singular and nominative/accusative plural.}{Singular nominative and accusative: \Latin{verbum}; plural nominative and accusative: \Latin{verba}.}
\Exercise{Transform \Latin{Dominus servum vocat} so that plural lords call plural servants; label both noun cases.}{\Latin{Domini servos vocant}. \Latin{Domini} is nominative plural; \Latin{servos} is accusative plural.}
\Exercise{Translate and parse the noun phrase \Latin{per omnia sæcula sæculorum}.}{\English{through all ages of ages} or idiomatically \English{forever and ever}. \Latin{per} governs neuter plural accusative \Latin{omnia sæcula}; \Latin{sæculorum} is neuter plural genitive.}
\PassageExercise{In principio creavit Deus cælum et terram.}{Parse each noun and translate.}{\Latin{principio}: neuter singular ablative after \Latin{in}; \Latin{Deus}: masculine singular nominative subject; \Latin{cælum}: neuter singular accusative object; \Latin{terram}: feminine singular accusative coordinated object. \English{In the beginning God created heaven and earth.}}
\end{worksheet}

\begin{worksheet}{I.10}{B}{Unit 3: First and Second Declensions}{Variant B}{Declensional production and ending-based correction}{Use macron-free course spelling and explain corrections from the paradigm.}
\Exercise{Decline \Latin{gratia, gratiæ f.} in the singular, in order nominative, genitive, dative, accusative, ablative, vocative.}{\Latin{gratia, gratiæ, gratiæ, gratiam, gratia, gratia}.}
\Exercise{Decline \Latin{Deus, Dei m.} in the plural in the same case order; give accepted alternatives where relevant.}{\Latin{dei} (also \Latin{dii}), \Latin{deorum} (also \Latin{deum}), \Latin{deis} (also \Latin{diis}), \Latin{deos}, \Latin{deis} (also \Latin{diis}), \Latin{dei} (also \Latin{dii}).}
\Exercise{Give the singular vocative of \Latin{filius, filii m.}, \Latin{servus, servi m.}, and \Latin{Deus, Dei m.}.}{\Latin{fili, serve, Deus}. Nouns in \Latin{-ius} normally contract to \Latin{-i}; regular \Latin{-us} becomes \Latin{-e}; liturgical \Latin{Deus} keeps \Latin{Deus}.}
\Exercise{Change \Latin{verbum bonum} to genitive singular, dative plural, and accusative plural.}{\Latin{verbi boni}; \Latin{verbis bonis}; \Latin{verba bona}.}
\Exercise{Correct every error in \Latin{in principium erat verba}, intended \English{in the beginning was the Word}.}{\Latin{In principio erat Verbum}. Stationary/metaphorical \Latin{in} takes ablative \Latin{principio}; singular subject is \Latin{Verbum}.}
\PassageExercise{Deo gratias.}{List every possible isolated analysis of \Latin{Deo}, choose the contextual one, identify \Latin{gratias}, and translate.}{\Latin{Deo} can be dative or ablative singular; here it is dative recipient. \Latin{gratias} is first-declension feminine accusative plural. \English{Thanks be to God.}}
\end{worksheet}

\begin{worksheet}{I.11}{C}{Unit 3: First and Second Declensions}{Variant C}{Mixed first/second-declension parsing and transformation}{Separate stem and ending when asked.}
\Exercise{For \Latin{ecclesiæ, animis, filiorum, dona}, list all possible case-number analyses before context.}{\Latin{ecclesiæ}: genitive or dative singular, or nominative/vocative plural. \Latin{animis}: dative or ablative plural. \Latin{filiorum}: genitive plural. \Latin{dona}: neuter nominative/accusative/vocative plural, or singular imperative of \Latin{dono} in another context.}
\Exercise{Separate stem and ending in \Latin{animarum, domino, verbis, sæcula}; state declension.}{\Latin{anim-arum}, first; \Latin{domin-o}, second masculine; \Latin{verb-is}, second neuter; \Latin{sæcul-a}, second neuter.}
\Exercise{Translate \English{the words of the Lord} and \English{to the souls of the servants}.}{\Latin{verba Domini}; \Latin{animis servorum}.}
\Exercise{Transform \Latin{Maria filium laudat}, \English{Mary praises the son}, into \English{The sons praise Mary}.}{\Latin{Filii Mariam laudant}.}
\Exercise{Diagnose \Latin{Domine Deus, dona anima pacem}, intended \English{Lord God, grant peace to the soul}.}{\Latin{Domine Deus, dona animæ pacem}. \Latin{Domine Deus} is direct address; recipient \Latin{animæ} must be dative; \Latin{pacem} remains accusative.}
\PassageExercise{Per omnia sæcula sæculorum. Amen.}{Parse all declined forms, identify the indeclinable item, and give close and idiomatic translations.}{\Latin{omnia sæcula}: neuter plural accusative after \Latin{per}; \Latin{sæculorum}: neuter plural genitive. \Latin{Amen} is indeclinable. Close: \English{through all ages of ages}; idiomatic: \English{forever and ever. Amen.}}
\end{worksheet}

\begin{worksheet}{I.12}{D}{Unit 3: First and Second Declensions}{Variant D}{Paradigm recall, neuter control, and syntax}{This may replace any other Unit 3 variant.}
\Exercise{Write the full plural paradigm of \Latin{missa, missæ f.} in case order.}{\Latin{missæ, missarum, missis, missas, missis, missæ}.}
\Exercise{Write the full singular paradigm of \Latin{verbum, verbi n.} in case order.}{\Latin{verbum, verbi, verbo, verbum, verbo, verbum}.}
\Exercise{Give the requested form: \English{of God}; \English{O Lord}; \English{to the son}; \English{with gifts}; \English{the churches} as direct object.}{\Latin{Dei; Domine; filio; cum donis; ecclesias}.}
\Exercise{Explain why \Latin{dona} in \Latin{dona nobis pacem} and \Latin{dona Dei} should not receive the same analysis.}{In \Latin{dona nobis}, \Latin{dona} is the singular imperative \English{grant}. In \Latin{dona Dei}, it is neuter plural nominative or accusative noun \English{gifts}, modified by genitive \Latin{Dei}. Syntax and dictionary possibilities decide.}
\Exercise{Correct \Latin{Filius custodit ecclesia Domini}, intended \English{The Son guards the Church of the Lord}.}{\Latin{Filius custodit Ecclesiam Domini}. Direct object \Latin{Ecclesiam} must be accusative; \Latin{Domini} is correctly genitive.}
\PassageExercise{Peccata mundi tollis.}{This is a controlled transformation of \Latin{qui tollis peccata mundi}, not a received quotation; apply the neuter rule and parse the genitive.}{\Latin{peccata}: neuter plural accusative object; neuter plural nominative would look identical, but second-person verb and context make it object. \Latin{mundi}: masculine singular genitive. Understood subject: \Latin{tu}. \English{You take away the sins of the world.}}
\end{worksheet}

\end{courseunit}

\begin{courseunit}{I.4}
\chapter{Unit 4: Third Declension}

\begin{worksheet}{I.13}{A}{Unit 4: Third Declension}{Variant A}{Stem recovery, consonant stems, i-stems, and liturgical syntax}{Derive each stem from the genitive before choosing an ending.}
\Exercise{Derive the stem and state gender for \Latin{rex, regis m.; pax, pacis f.; nomen, nominis n.; corpus, corporis n.; altare, altaris n.; lux, lucis f.}.}{\Latin{reg-}, masculine; \Latin{pac-}, feminine; \Latin{nomin-}, neuter; \Latin{corpor-}, neuter; \Latin{altar-}, neuter; \Latin{luc-}, feminine.}
\Exercise{Give accusative singular, genitive plural, and dative/ablative plural of \Latin{rex, regis m.}.}{\Latin{regem, regum, regibus}.}
\Exercise{Give nominative/accusative plural and genitive plural of consonant-stem \Latin{nomen} and i-stem \Latin{altare}; state the contrasting rule.}{\Latin{nomina, nominum}; \Latin{altaria, altarium}. Neuter consonant stems use \Latin{-a/-um}; neuter i-stems use \Latin{-ia/-ium}.}
\Exercise{Translate and parse \Latin{ad altare Dei}.}{\Latin{ad} governs neuter singular accusative \Latin{altare}; \Latin{Dei} is masculine singular genitive. \English{to the altar of God}.}
\Exercise{Correct \Latin{in nominem Patris}, intended \English{in the name of the Father}.}{\Latin{in nomine Patris}. Here \Latin{in} denotes location or state and governs ablative \Latin{nomine}; \Latin{Patris} is genitive singular.}
\PassageExercise{Pax hominibus bonæ voluntatis.}{Recover every third-declension dictionary form, parse the printed forms, and translate.}{\Latin{pax, pacis f.}: nominative singular; \Latin{homo, hominis m.}: \Latin{hominibus}, dative plural; \Latin{voluntas, voluntatis f.}: \Latin{voluntatis}, genitive singular. \Latin{bonæ} agrees with \Latin{voluntatis}. \English{Peace [be] to people of good will.}}
\end{worksheet}

\begin{worksheet}{I.14}{B}{Unit 4: Third Declension}{Variant B}{Third-declension paradigms and neuter control}{Use dictionary forms, not resemblance to a first- or second-declension ending.}
\Exercise{Decline \Latin{nomen, nominis n.} in the singular in case order.}{\Latin{nomen, nominis, nomini, nomen, nomine, nomen}.}
\Exercise{Decline \Latin{altare, altaris n.} in the plural in case order.}{\Latin{altaria, altarium, altaribus, altaria, altaribus, altaria}.}
\Exercise{Give the requested forms of \Latin{pater, patris m.}: genitive singular, accusative singular, ablative singular, nominative plural, genitive plural.}{\Latin{patris, patrem, patre, patres, patrum}.}
\Exercise{Translate \English{the bodies of the kings}, \English{to the fathers}, and \English{with perpetual light}.}{\Latin{corpora regum}; \Latin{patribus}; \Latin{cum luce perpetua}.}
\Exercise{Diagnose \Latin{corpus Christi custodiunt animam}, if singular Body is the subject.}{\Latin{Corpus Christi custodit animam}. Neuter singular \Latin{corpus} requires singular verb \Latin{custodit}; \Latin{Christi} is genitive and \Latin{animam} accusative.}
\PassageExercise{In nomine Patris, et Filii, et Spiritus Sancti.}{Identify the third-declension forms, then account for every case relation and translate.}{Third-declension \Latin{nomine} is neuter singular ablative after \Latin{in}; \Latin{Patris} is masculine singular genitive. \Latin{Filii} and \Latin{Spiritus Sancti} are coordinated genitives. \English{In the name of the Father, and of the Son, and of the Holy Spirit.}}
\end{worksheet}

\begin{worksheet}{I.15}{C}{Unit 4: Third Declension}{Variant C}{Ambiguous forms, i-stems, and transformations}{List possibilities before selecting the contextual analysis.}
\Exercise{For isolated \Latin{reges, nominibus, altaria, lucis}, list all possible case-number analyses.}{\Latin{reges}: nominative/accusative/vocative plural. \Latin{nominibus}: dative/ablative plural. \Latin{altaria}: nominative/accusative/vocative plural. \Latin{lucis}: genitive singular or nominative/accusative plural only if a different lexeme were involved; for \Latin{lux, lucis f.}, genitive singular.}
\Exercise{State the i-stem evidence and give ablative singular, nominative plural, and genitive plural for \Latin{altare, altaris n.}.}{Nominative in \Latin{-ar} marks a neuter i-stem. Forms: \Latin{altari, altaria, altarium}.}
\Exercise{Transform \Latin{Rex custodit civitatem} so that plural kings guard plural cities.}{\Latin{Reges custodiunt civitates}. \Latin{Reges} is nominative plural; \Latin{civitates} accusative plural.}
\Exercise{Give dictionary form, stem, and parse for \Latin{hominibus} in \Latin{dona hominibus pacem}.}{\Latin{homo, hominis m.}; stem \Latin{homin-}; dative plural recipient of \Latin{dona}.}
\Exercise{Correct \Latin{ad altaribus Dei} and explain the government.}{\Latin{ad altare Dei} for one altar, or \Latin{ad altaria Dei} for several. \Latin{Ad} governs accusative, never ablative \Latin{altaribus}.}
\PassageExercise{Corpus Domini nostri Iesu Christi custodiat animam meam in vitam æternam.}{Identify and parse every third-declension noun, classify the other nouns, and give a close translation.}{The only third-declension noun is \Latin{Corpus, corporis n.}, nominative singular subject. \Latin{Domini} and \Latin{Christi} are second-declension genitives; \Latin{Iesu} is the irregular genitive of \Latin{Iesus}; \Latin{animam} and \Latin{vitam} are first-declension accusatives. The genitive chain modifies \Latin{Corpus}; \Latin{animam meam} is the object; \Latin{in vitam æternam} is an accusative goal. \English{May the Body of our Lord Jesus Christ preserve my soul unto everlasting life.}}
\end{worksheet}

\begin{worksheet}{I.16}{D}{Unit 4: Third Declension}{Variant D}{Stem-first production and syntax-driven parsing}{Parallel practice with a different vocabulary set.}
\Exercise{From \Latin{homo, hominis m.; civitas, civitatis f.; mors, mortis f.; sacerdos, sacerdotis m.; sanguis, sanguinis m.}, derive each stem.}{\Latin{homin-, civitat-, mort-, sacerdot-, sanguin-}.}
\Exercise{Give accusative singular and plural, genitive plural, and ablative singular of \Latin{civitas, civitatis f.}.}{\Latin{civitatem, civitates, civitatum, civitate}.}
\Exercise{Apply the neuter rule to \Latin{corpus, corporis n.}: give accusative singular, nominative plural, accusative plural, and genitive plural.}{\Latin{corpus, corpora, corpora, corporum}.}
\Exercise{Translate \English{the name of the king}, \English{peace to the people}, and \English{at the altars}.}{\Latin{nomen regis}; \Latin{pax hominibus}; \Latin{in altaribus}.}
\Exercise{Correct the two ending errors in \Latin{Lux perpetuum luceat eis}.}{\Latin{Lux perpetua luceat eis}. \Latin{Lux} is feminine singular nominative, so \Latin{perpetua} agrees; \Latin{eis} is already the correct dative plural.}
\PassageExercise{Introibo ad altare Dei, ad Deum qui lætificat iuventutem meam.}{Parse \Latin{altare, Dei, Deum}, and \Latin{iuventutem}; resolve the case of \Latin{qui}; translate.}{\Latin{altare}: neuter singular accusative after \Latin{ad}; \Latin{Dei}: genitive singular; \Latin{Deum}: masculine singular accusative after \Latin{ad}; \Latin{qui}: masculine singular nominative subject of \Latin{lætificat}, with antecedent \Latin{Deum}; \Latin{iuventutem}: feminine singular accusative object. \English{I shall go unto the altar of God, to God who gives joy to my youth.}}
\end{worksheet}

\end{courseunit}

\begin{courseunit}{I.5}
\chapter{Unit 5: Fourth and Fifth Declensions}

\begin{worksheet}{I.17}{A}{Unit 5: Fourth and Fifth Declensions}{Variant A}{Fourth/fifth forms and five-declension discrimination}{Write gender as well as case and number.}
\Exercise{Give genitive singular, dative singular, accusative singular, ablative singular, and genitive plural of \Latin{spiritus, spiritus m.}.}{\Latin{spiritus, spiritui, spiritum, spiritu, spirituum}.}
\Exercise{Give nominative/accusative singular, nominative/accusative plural, genitive plural, and dative/ablative plural of \Latin{genu, genus n.}.}{\Latin{genu}; \Latin{genua}; \Latin{genuum}; \Latin{genibus}.}
\Exercise{Decline \Latin{res, rei f.} in the singular in case order.}{\Latin{res, rei, rei, rem, re, res}.}
\Exercise{Parse and translate \Latin{cum spiritu tuo}; explain why \Latin{spiritum} would be wrong.}{\Latin{spiritu tuo} is masculine singular ablative after \Latin{cum}: \English{with your spirit}. \Latin{spiritum} is accusative, which \Latin{cum} does not govern.}
\Exercise{Correct \Latin{hæc est diem quam fecit Dominus}.}{\Latin{Hæc est dies quam fecit Dominus}. Predicate noun \Latin{dies} must be nominative after \Latin{est}; \Latin{quam} remains accusative object of \Latin{fecit}.}
\PassageExercise{Pater, in manus tuas commendo spiritum meum.}{Parse the fourth-declension nouns, explain both adjective forms, and translate.}{\Latin{manus}: feminine plural accusative after directional \Latin{in}; \Latin{tuas} agrees. \Latin{spiritum}: masculine singular accusative object; \Latin{meum} agrees. \English{Father, into your hands I commend my spirit.}}
\end{worksheet}

\begin{worksheet}{I.18}{B}{Unit 5: Fourth and Fifth Declensions}{Variant B}{Ambiguous fourth-declension endings and fifth-declension production}{Use syntax to choose among identical-looking forms.}
\Exercise{List every possible analysis of isolated \Latin{spiritus} within the paradigm of \Latin{spiritus, spiritus m.}.}{Nominative singular, genitive singular, nominative plural, accusative plural, vocative singular, or vocative plural.}
\Exercise{Resolve \Latin{spiritus} in (a) \Latin{Spiritus Sanctus venit}, (b) \Latin{dona Spiritus Sancti}, and (c) \Latin{invocat spiritus}.}{(a) nominative singular subject; (b) genitive singular modifying \Latin{dona}; (c) accusative plural object, \English{he calls the spirits}.}
\Exercise{Give plural nominative, genitive, dative, accusative, and ablative of \Latin{res, rei f.}.}{\Latin{res, rerum, rebus, res, rebus}.}
\Exercise{Give dative/ablative plural of \Latin{dies, diei}; \Latin{manus, manus}; and \Latin{genu, genus}.}{\Latin{diebus, manibus, genibus}.}
\Exercise{Diagnose \Latin{omne genus flectantur}, intended \English{every knee should bend}.}{\Latin{Omne genu flectatur}. The noun is fourth-declension neuter \Latin{genu, genus}; singular subject takes singular verb. \Latin{Genus} would be the genitive \English{of a knee} here, not the nominative.}
\PassageExercise{Hæc est dies quam fecit Dominus: exsultemus et lætemur in ea.}{Parse both forms referring to \Latin{dies}, identify every governed prepositional form, and translate.}{\Latin{hæc}: feminine singular nominative agreeing with contextual feminine \Latin{dies}; \Latin{quam}: feminine singular accusative object of \Latin{fecit}; \Latin{in ea}: feminine singular ablative after \Latin{in} denoting the situation or time in which. \English{This is the day which the Lord has made; let us rejoice and be glad in it.}}
\end{worksheet}

\begin{worksheet}{I.19}{C}{Unit 5: Fourth and Fifth Declensions}{Variant C}{Five-declension comparison and controlled composition}{Name the declension from the dictionary genitive, not the nominative.}
\Exercise{Classify by declension and gender: \Latin{anima, animæ f.; dominus, domini m.; pax, pacis f.; manus, manus f.; dies, diei m.}.}{First feminine; second masculine; third feminine; fourth feminine; fifth masculine.}
\Exercise{Give the genitive singular and nominative plural of each noun in the preceding item.}{\Latin{animæ/animæ}; \Latin{domini/domini}; \Latin{pacis/paces}; \Latin{manus/manus}; \Latin{diei/dies}.}
\Exercise{Translate \English{of the Holy Spirit}, \English{with both hands}, \English{the affairs of the Church}, and \English{on that day}.}{\Latin{Spiritus Sancti}; \Latin{ambabus manibus}; \Latin{res Ecclesiæ}; \Latin{in eo die}.}
\Exercise{Transform \Latin{omne genu} to plural and \Latin{hæc res} to accusative singular and plural.}{\Latin{omnia genua}; \Latin{hanc rem}; \Latin{has res}.}
\Exercise{Correct \Latin{in manibus tuas spiritum meum commendo}, preserving the meaning \English{into your hands I commend my spirit}.}{\Latin{In manus tuas spiritum meum commendo}. Motion toward requires accusative \Latin{manus tuas}; \Latin{spiritum meum} is already accusative object.}
\PassageExercise{Et cum spiritu tuo.}{List all possible analyses of \Latin{spiritu} before selecting one, connect the adjective, and translate.}{Within this paradigm \Latin{spiritu} is ablative singular. \Latin{Cum} confirms that analysis; \Latin{tuo} is masculine singular ablative in agreement. \English{And with your spirit.}}
\end{worksheet}

\begin{worksheet}{I.20}{D}{Unit 5: Fourth and Fifth Declensions}{Variant D}{Rare declensions in high-frequency liturgical phrases}{Treat every form as part of the complete noun system.}
\Exercise{Write the full singular and plural paradigms of \Latin{genu, genus n.}.}{Singular: \Latin{genu, genus, genu, genu, genu, genu}. Plural: \Latin{genua, genuum, genibus, genua, genibus, genua}.}
\Exercise{Write genitive singular, dative singular, accusative singular, ablative singular, and genitive plural of \Latin{dies, diei m.}.}{\Latin{diei, diei, diem, die, dierum}.}
\Exercise{Explain how context distinguishes \Latin{res} in \Latin{res bona est}, \Latin{videt res bonas}, and \Latin{res bonæ sunt}.}{First \Latin{res} is nominative singular; second is accusative plural; third is nominative plural. Verb number and adjective endings resolve the identical noun form.}
\Exercise{Translate \English{the spirit of truth}, \English{to the Holy Spirit}, and \English{the knees of the servants}.}{\Latin{spiritus veritatis}; \Latin{Spiritui Sancto}; \Latin{genua servorum}.}
\Exercise{Diagnose \Latin{cum Spiritum Sanctum} and \Latin{rei bonus}; correct both.}{\Latin{cum Spiritu Sancto}, ablative after \Latin{cum}; \Latin{rei bonæ}, feminine singular genitive or dative agreement.}
\PassageExercise{Ut in nomine Iesu omne genu flectatur.}{Parse \Latin{nomine, Iesu, omne genu}, and the finite verb; translate.}{\Latin{nomine}: neuter singular ablative after \Latin{in}; \Latin{Iesu}: genitive singular; \Latin{omne genu}: neuter singular nominative subject; \Latin{flectatur}: third-person singular present passive subjunctive. \English{That at the name of Jesus every knee should bend.}}
\end{worksheet}

\end{courseunit}

\begin{courseunit}{I.6}
\chapter{Unit 6: Adjectives and Agreement}

\begin{worksheet}{I.21}{A}{Unit 6: Adjectives and Agreement}{Variant A}{First/second and third-declension adjective systems}{Always state the noun that controls agreement.}
\Exercise{Give all three genders of \Latin{sanctus} in nominative singular, accusative singular, and nominative plural.}{Nominative singular: \Latin{sanctus, sancta, sanctum}; accusative singular: \Latin{sanctum, sanctam, sanctum}; nominative plural: \Latin{sancti, sanctæ, sancta}.}
\Exercise{Give masculine/feminine and neuter forms of \Latin{omnis, omne} in accusative singular, nominative plural, and genitive plural.}{Accusative singular: \Latin{omnem/omne}; nominative plural: \Latin{omnes/omnia}; genitive plural: \Latin{omnium} for all genders.}
\Exercise{Make \Latin{æternus, -a, -um} agree with \Latin{requiem}, \Latin{vita}, \Latin{lumen}, and \Latin{sæcula} in the cases shown by those isolated forms. List ambiguity where it remains.}{\Latin{requiem æternam}, accusative singular; \Latin{vita æterna}, nominative or ablative singular; \Latin{lumen æternum}, nominative or accusative singular; \Latin{sæcula æterna}, nominative or accusative plural.}
\Exercise{Parse \Latin{Patris omnipotentis}: explain why the endings differ but agreement is exact.}{Both are masculine singular genitive. \Latin{Patris} is a third-declension noun and \Latin{omnipotentis} a third-declension adjective; agreement matches coordinates, not letters alone.}
\Exercise{Correct \Latin{lux perpetuum} and \Latin{omnia homines}; state each controlling noun's gender.}{\Latin{lux perpetua}, because \Latin{lux} is feminine; \Latin{omnes homines}, because \Latin{homines} is masculine plural nominative or accusative.}
\PassageExercise{Requiem æternam dona eis, Domine, et lux perpetua luceat eis.}{Parse both adjective-noun pairs, identify the vocative, and translate.}{\Latin{requiem æternam}: feminine singular accusative object; \Latin{lux perpetua}: feminine singular nominative subject; \Latin{Domine}: masculine singular vocative. \English{Grant them eternal rest, O Lord, and may perpetual light shine upon them.}}
\end{worksheet}

\begin{worksheet}{I.22}{B}{Unit 6: Adjectives and Agreement}{Variant B}{Agreement across unlike declensions and predicate adjectives}{This variant checks the same adjective systems with different forms.}
\Exercise{Decline \Latin{Spiritus Sanctus} in singular genitive, dative, accusative, and ablative.}{\Latin{Spiritus Sancti, Spiritui Sancto, Spiritum Sanctum, Spiritu Sancto}.}
\Exercise{Give these forms of one-termination \Latin{felix, felicis}: masculine/feminine accusative singular, neuter nominative singular, ablative singular all genders, neuter nominative plural, genitive plural.}{\Latin{felicem, felix, felici, felicia, felicium}.}
\Exercise{Transform \Latin{anima sancta} to accusative plural; \Latin{nomen sanctum} to genitive plural; \Latin{rex sanctus} to dative plural.}{\Latin{animas sanctas}; \Latin{nominum sanctorum}; \Latin{regibus sanctis}.}
\Exercise{In \Latin{Dignum et iustum est}, classify the adjective use and explain the neuter singular.}{They are predicate adjectives after \Latin{est}. Neuter singular refers to an understood action or \Latin{hoc}, \English{this/it}: \English{It is fitting and just.}}
\Exercise{Diagnose \Latin{Omnipotente sempiterne Deus} as a direct address.}{\Latin{Omnipotens sempiterne Deus}. The one-termination adjective's masculine singular vocative equals nominative \Latin{omnipotens}; \Latin{sempiterne} is the regular vocative.}
\PassageExercise{Sanctus, Sanctus, Sanctus Dominus Deus Sabaoth.}{Parse every adjective and noun that declines; identify the indeclinable item; translate.}{Each \Latin{Sanctus} is masculine singular nominative in acclamatory/predicate agreement with \Latin{Dominus}; \Latin{Deus} is masculine singular nominative in apposition. \Latin{Sabaoth} is indeclinable. \English{Holy, Holy, Holy Lord God of hosts.}}
\end{worksheet}

\begin{worksheet}{I.23}{C}{Unit 6: Adjectives and Agreement}{Variant C}{Attributive, predicate, and substantival adjectives}{Name both morphology and use.}
\Exercise{Classify the adjective as attributive, predicate, or substantival: \Latin{vita æterna}; \Latin{pleni sunt cæli}; \Latin{sancti laudant}; \Latin{homines bonæ voluntatis}.}{\Latin{æterna}: attributive; \Latin{pleni}: predicate; \Latin{sancti}: substantival, \English{the saints/holy ones}; \Latin{bonæ}: attributive to \Latin{voluntatis}.}
\Exercise{Supply \Latin{omnis, omne}: \English{every name}, \English{of every name}, \English{all souls}, \English{with all spirits}.}{\Latin{omne nomen}; \Latin{omnis nominis}; \Latin{omnes animæ}; \Latin{cum omnibus spiritibus}.}
\Exercise{Change \Latin{Corpus Domini nostri} to plural \English{Bodies of our Lords}; parse the possessive and the nouns it connects.}{\Latin{Corpora Dominorum nostrorum}. \Latin{Corpora} is neuter plural nominative or accusative; \Latin{Dominorum} is masculine plural genitive; possessive \Latin{nostrorum} agrees with \Latin{Dominorum}, not directly with \Latin{Corpora}.}
\Exercise{Translate \English{all holy names}, \English{the happy souls}, and \English{of almighty God}.}{\Latin{omnia nomina sancta}; \Latin{animæ felices}; \Latin{Dei omnipotentis}.}
\Exercise{Correct \Latin{plena sunt cæli et terra gloria tuum}.}{\Latin{Pleni sunt cæli et terra gloria tua}. Predicate adjective for the mixed compound subject is masculine plural \Latin{pleni}; \Latin{gloria tua} is feminine singular ablative.}
\PassageExercise{Pleni sunt cæli et terra gloria tua.}{Identify subject, linking verb, predicate adjective, and the remaining ablative phrase; give close English.}{Compound subject: \Latin{cæli et terra}; linking verb: \Latin{sunt}; predicate adjective: masculine plural nominative \Latin{pleni}; ablative complement: \Latin{gloria tua}. \English{Heaven and earth are full of your glory.}}
\end{worksheet}

\begin{worksheet}{I.24}{D}{Unit 6: Adjectives and Agreement}{Variant D}{Adjective paradigms, agreement chains, and correction}{Produce a complete rationale for every changed ending.}
\Exercise{Write the full singular masculine paradigm of \Latin{sanctus} and the full singular feminine paradigm.}{Masculine: \Latin{sanctus, sancti, sancto, sanctum, sancto, sancte}. Feminine: \Latin{sancta, sanctæ, sanctæ, sanctam, sancta, sancta}.}
\Exercise{Write the plural paradigm of \Latin{omnis, omne}, giving masculine/feminine then neuter in every case.}{Nom. \Latin{omnes/omnia}; gen. \Latin{omnium/omnium}; dat. \Latin{omnibus/omnibus}; acc. \Latin{omnes/omnia}; abl. \Latin{omnibus/omnibus}; voc. \Latin{omnes/omnia}.}
\Exercise{Parse every adjective in \Latin{Corpus Domini nostri Iesu Christi custodiat animam meam in vitam æternam}.}{\Latin{nostri}: masculine singular genitive with \Latin{Domini}; \Latin{meam}: feminine singular accusative with \Latin{animam}; \Latin{æternam}: feminine singular accusative with \Latin{vitam}.}
\Exercise{Transform \Latin{lux perpetua} to genitive plural, \Latin{spiritus sanctus} to accusative plural, and \Latin{omne genu} to nominative plural.}{\Latin{lucum perpetuarum}; \Latin{spiritus sanctos}; \Latin{omnia genua}.}
\Exercise{Diagnose \Latin{animas felicis custodit Deus bonus}.}{\Latin{Animas felices custodit Deus bonus}. \Latin{felices} must be feminine plural accusative with \Latin{animas}; \Latin{bonus} correctly agrees with nominative singular \Latin{Deus}.}
\PassageExercise{Omnipotens sempiterne Deus, miserere nobis.}{Treat this as a project-composed syntax example; parse both adjectives and translate.}{\Latin{Deus} is masculine singular direct address with nominative-shaped vocative. \Latin{omnipotens} is a one-termination third-declension adjective whose vocative equals nominative; \Latin{sempiterne} is first/second-declension masculine singular vocative. \English{Almighty, eternal God, have mercy on us.}}
\end{worksheet}

\end{courseunit}

\begin{courseunit}{I.7}
\chapter{Unit 7: \Latin{Sum} and \Latin{Possum}}

\begin{worksheet}{I.25}{A}{Unit 7: \Latin{Sum} and \Latin{Possum}}{Variant A}{Present-system forms, predicates, and complementary infinitives}{Parse every finite verb by person, number, tense, voice, and mood.}
\Exercise{Write the present, imperfect, and future of \Latin{sum} in the first-person singular, second-person plural, and third-person plural.}{Present: \Latin{sum, estis, sunt}. Imperfect: \Latin{eram, eratis, erant}. Future: \Latin{ero, eritis, erunt}. All are active indicative forms.}
\Exercise{Parse \Latin{sum, eras, erimus, sunt}; give one unambiguous English rendering for each.}{\Latin{sum}: 1 singular present, \English{I am}; \Latin{eras}: 2 singular imperfect, \English{you were}; \Latin{erimus}: 1 plural future, \English{we shall be}; \Latin{sunt}: 3 plural present, \English{they are}.}
\Exercise{Change \Latin{sum dignus} through all six present persons, using a masculine or mixed group in the plural.}{\Latin{sum dignus, es dignus, est dignus, sumus digni, estis digni, sunt digni}. A female speaker/group would use the corresponding feminine adjective.}
\Exercise{Correct \Latin{Maria est bonam} and explain why a linking verb does not license the accusative.}{\Latin{Maria est bona}. Predicate adjective \Latin{bona} is feminine singular nominative agreeing with subject \Latin{Maria}; \Latin{sum} links rather than takes a direct object.}
\Exercise{Translate \English{you (plural) can praise God} and \English{we were able to hear the word}.}{\Latin{Potestis Deum laudare}; \Latin{Poteramus verbum audire}. \Latin{Laudare} and \Latin{audire} are complementary infinitives.}
\PassageExercise{Dignum et iustum est.}{Parse the verb, identify the adjective use and understood reference, and translate.}{\Latin{est}: third-person singular present active indicative. \Latin{dignum} and \Latin{iustum}: neuter singular nominative predicate adjectives referring to an understood \Latin{hoc} or action. \English{It is fitting and just.}}
\end{worksheet}

\begin{worksheet}{I.26}{B}{Unit 7: \Latin{Sum} and \Latin{Possum}}{Variant B}{Perfect-system forms and linking-clause syntax}{This is equivalent to Variant A but ranges across all six indicative tenses.}
\Exercise{Write all six persons of the perfect active indicative of \Latin{sum}.}{\Latin{fui, fuisti, fuit, fuimus, fuistis, fuerunt}.}
\Exercise{Parse \Latin{fueramus, fuero, fuerint, fuit}; note any form that can be ambiguous outside this unit's indicative frame.}{\Latin{fueramus}: 1 plural pluperfect active indicative; \Latin{fuero}: 1 singular future perfect active indicative; \Latin{fuerint}: 3 plural future perfect active indicative, also potentially perfect subjunctive; \Latin{fuit}: 3 singular perfect active indicative.}
\Exercise{Transform \Latin{sancti erunt}, \English{they will be holy}, into imperfect, perfect, and pluperfect.}{\Latin{sancti erant; sancti fuerunt; sancti fuerant}.}
\Exercise{Give the present and perfect infinitives and the singular/plural imperatives of \Latin{sum}.}{Present infinitive \Latin{esse}; perfect infinitive \Latin{fuisse}; imperatives \Latin{es, este}.}
\Exercise{Diagnose \Latin{Via et veritas est bonæ}, intended \English{The way and truth are good}.}{\Latin{Via et veritas sunt bonæ}. Compound subject requires plural \Latin{sunt}; predicate adjective is feminine plural nominative \Latin{bonæ}.}
\PassageExercise{Ego sum via, et veritas, et vita.}{Identify subject, linking verb, and all predicate nouns; explain their case and translate.}{\Latin{ego}: nominative subject; \Latin{sum}: 1 singular present active indicative; \Latin{via, veritas, vita}: coordinated nominative predicate nouns. \English{I am the way, and the truth, and the life.}}
\end{worksheet}

\begin{worksheet}{I.27}{C}{Unit 7: \Latin{Sum} and \Latin{Possum}}{Variant C}{The compound \Latin{possum}, tense contrasts, and ellipsis}{Supply an infinitive wherever ability requires one.}
\Exercise{Write all six present forms of \Latin{possum}.}{\Latin{possum, potes, potest, possumus, potestis, possunt}.}
\Exercise{Give the requested forms: 3 plural imperfect, 2 singular future, 1 plural perfect, 3 singular pluperfect, 2 plural future perfect of \Latin{possum}.}{\Latin{poterant, poteris, potuimus, potuerat, potueritis}.}
\Exercise{Translate and distinguish \Latin{poterat facere}, \Latin{potuit facere}, and \Latin{potuerat facere}.}{\English{He was able/could do}; \English{he was able/managed to do}; \English{he had been able to do}. They are imperfect, perfect, and pluperfect respectively, all with complementary \Latin{facere}.}
\Exercise{Complete with the correct form of \Latin{possum}:\par
\Latin{ego \Blank audire; nos \Blank laudare; vos cras \Blank venire}.}{\Latin{ego possum audire; nos possumus laudare; vos cras poteritis venire}.}
\Exercise{Correct \Latin{Possunt Deum laudant}.}{\Latin{Possunt Deum laudare}. A finite form of \Latin{possum} takes a complementary infinitive, not another finite present verb.}
\PassageExercise{Sine me nihil potestis facere.}{Parse \Latin{potestis}, identify its infinitive and the roles of \Latin{me} and \Latin{nihil}, then translate.}{\Latin{potestis}: 2 plural present active indicative of \Latin{possum}; complement \Latin{facere}; \Latin{me}: ablative after \Latin{sine}; \Latin{nihil}: neuter accusative object of \Latin{facere}. \English{Without me you can do nothing.}}
\end{worksheet}

\begin{worksheet}{I.28}{D}{Unit 7: \Latin{Sum} and \Latin{Possum}}{Variant D}{Copular ellipsis and cumulative conjugation}{Bracket supplied forms and keep printed forms distinct.}
\Exercise{Give the conventional optative construal of \Latin{Dominus vobiscum} and distinguish interpretation from printed form.}{Analytic paraphrase: \Latin{Dominus [sit] vobiscum}. \Latin{Sit} would be 3 singular present subjunctive expressing a wish, but it is not printed or uniquely forced: \English{The Lord be with you}.}
\Exercise{Write second-person singular and third-person plural forms of \Latin{sum} in all six indicative tenses.}{2 singular: \Latin{es, eras, eris, fuisti, fueras, fueris}. 3 plural: \Latin{sunt, erant, erunt, fuerunt, fuerant, fuerint}.}
\Exercise{Parse and translate \Latin{pleni sunt}, \Latin{parati eramus}, \Latin{sancta fuerat}, and \Latin{beati erunt}.}{\Latin{sunt}: 3 plural present, \English{they are full}; \Latin{eramus}: 1 plural imperfect, \English{we were ready}; \Latin{fuerat}: 3 singular pluperfect, \English{it/she had been holy}; \Latin{erunt}: 3 plural future, \English{they will be blessed}.}
\Exercise{Translate \English{I shall be able to come}, \English{they had been able to see}, and \English{be holy!} to several men.}{\Latin{Potero venire}; \Latin{potuerant videre}; \Latin{este sancti!}}
\Exercise{Diagnose \Latin{Cæli et terra plenum est}.}{\Latin{Cæli et terra pleni sunt}. Compound plural subject requires \Latin{sunt}; mixed-gender predicate is masculine plural \Latin{pleni}.}
\PassageExercise{Non sum dignus ut intres sub tectum meum.}{Parse the linking clause fully, identify the subjunctive without producing its full paradigm, and translate.}{Understood \Latin{ego} is subject; \Latin{sum} is 1 singular present active indicative; \Latin{dignus} is masculine singular nominative predicate adjective. \Latin{intres} is 2 singular present active subjunctive in an \Latin{ut}-clause completing \Latin{dignus}. \English{I am not worthy that you should enter under my roof.}}
\end{worksheet}

\end{courseunit}

\begin{courseunit}{I.8}
\chapter{Unit 8: The Present Active System}

\begin{worksheet}{I.29}{A}{Unit 8: The Present Active System}{Variant A}{Conjugation, present forms, and clause parsing}{Use all four principal parts when a dictionary entry is supplied.}
\Exercise{Classify by conjugation: \Latin{laudo, laudare; moneo, monere; rego, regere; capio, capere; audio, audire}.}{First; second; third; third-\Latin{io}; fourth. A marked dictionary distinguishes second \Latin{-ere} from third \Latin{-ere}; the first principal part also helps.}
\Exercise{Write all six present active indicative forms of \Latin{rego}.}{\Latin{rego, regis, regit, regimus, regitis, regunt}.}
\Exercise{Parse \Latin{laudamus, monet, capiunt, auditis}, giving lexical verb and translation.}{\Latin{laudamus}: 1 plural present active indicative, \English{we praise}; \Latin{monet}: 3 singular, \English{he/she warns}; \Latin{capiunt}: 3 plural, \English{they take}; \Latin{auditis}: 2 plural, \English{you hear}.}
\Exercise{Transform \Latin{Dominus animam custodit} to first-person plural subject, then to third-person plural subject.}{\Latin{Animam custodimus}; \Latin{Domini animam custodiunt} or, with a pronoun, \Latin{Ei animam custodiunt}. The verb endings must change.}
\Exercise{Correct \Latin{Nos Deum laudat et glorificatis}.}{\Latin{Nos Deum laudamus et glorificamus}. Both verbs share first-person plural subject \Latin{nos}.}
\PassageExercise{Laudamus te. Benedicimus te. Adoramus te. Glorificamus te.}{Parse every finite verb, identify the repeated object, and translate.}{All four verbs are 1 plural present active indicative; \Latin{te} is repeated singular accusative object. \English{We praise you. We bless you. We adore you. We glorify you.}}
\end{worksheet}

\begin{worksheet}{I.30}{B}{Unit 8: The Present Active System}{Variant B}{Imperfect and future active forms}{Contrast tense signs before translating.}
\Exercise{Give first-person singular and third-person plural imperfect active forms of \Latin{laudo, moneo, rego, capio}, and \Latin{audio}.}{\Latin{laudabam/laudabant; monebam/monebant; regebam/regebant; capiebam/capiebant; audiebam/audiebant}.}
\Exercise{Give the same persons in the future active.}{\Latin{laudabo/laudabunt; monebo/monebunt; regam/regent; capiam/capient; audiam/audient}.}
\Exercise{Parse \Latin{laudabatis, regebat, capietis, audient}.}{\Latin{laudabatis}: 2 plural imperfect active indicative; \Latin{regebat}: 3 singular imperfect; \Latin{capietis}: 2 plural future; \Latin{audient}: 3 plural future.}
\Exercise{Transform \Latin{verbum audimus} into imperfect and future, keeping person and number.}{\Latin{verbum audiebamus}; \Latin{verbum audiemus}.}
\Exercise{Diagnose \Latin{cras Dominum laudabant}, intended \English{tomorrow they will praise the Lord}.}{\Latin{Cras Dominum laudabunt}. \Latin{Laudabant} is imperfect, \English{they were praising}; future is \Latin{laudabunt}.}
\PassageExercise{Ego veniam et curabo eum.}{Parse both verbs, explain how coordination resolves possible ambiguity, identify the object, and translate.}{\Latin{veniam}: 1 singular future active indicative of \Latin{venio}; by itself it can resemble a present subjunctive, but future \Latin{curabo} and context support future. \Latin{curabo}: 1 singular future active indicative; \Latin{eum}: accusative object. \English{I will come and heal him.}}
\end{worksheet}

\begin{worksheet}{I.31}{C}{Unit 8: The Present Active System}{Variant C}{Infinitives, imperatives, and present-system transformations}{Distinguish a command from a finite statement.}
\Exercise{Give present active infinitive, singular imperative, and plural imperative of \Latin{laudo, moneo, rego, capio}, and \Latin{audio}.}{\Latin{laudare/lauda/laudate; monere/mone/monete; regere/rege/regite; capere/cape/capite; audire/audi/audite}.}
\Exercise{Translate to Latin: \English{praise the Lord!} to one and several; \English{hear the word!} to one and several.}{\Latin{Lauda Dominum! Laudate Dominum! Audi verbum! Audite verbum!}}
\Exercise{Change \Latin{laudatis} to a plural command, a plural imperfect statement, and a plural future statement without changing person.}{\Latin{laudate; laudabatis; laudabitis}.}
\Exercise{Parse \Latin{petite} and \Latin{accipietis}; explain why final \Latin{-ite} in the first and \Latin{-ietis} in the second do not mark the same category.}{\Latin{petite}: 2 plural present active imperative of \Latin{peto}. \Latin{accipietis}: 2 plural future active indicative of \Latin{accipio}; \Latin{-ie-} belongs to its future formation and \Latin{-tis} is the personal ending.}
\Exercise{Correct \Latin{Potes audi verbum}, intended \English{You can hear the word}.}{\Latin{Potes audire verbum}. \Latin{Potes} requires complementary infinitive \Latin{audire}, not imperative \Latin{audi}.}
\PassageExercise{Petite, et accipietis; quærite, et invenietis.}{Parse all four verbs and translate the two balanced commands and promises.}{\Latin{petite, quærite}: 2 plural present active imperatives. \Latin{accipietis}: 2 plural future active indicative of \Latin{accipio}; \Latin{invenietis}: 2 plural future active indicative of \Latin{invenio}. \English{Ask, and you will receive; seek, and you will find.}}
\end{worksheet}

\begin{worksheet}{I.32}{D}{Unit 8: The Present Active System}{Variant D}{Mixed present-system recognition in liturgical clauses}{Parse before using a memorized English text.}
\Exercise{From \Latin{credo, credere, credidi, creditum}, identify present stem, conjugation, and the form \Latin{credimus}.}{Present stem \Latin{cred-}; third conjugation; \Latin{credimus} is 1 plural present active indicative, \English{we believe}.}
\Exercise{Give 2 singular, 1 plural, and 3 plural in present, imperfect, and future for \Latin{audio}.}{Present: \Latin{audis, audimus, audiunt}. Imperfect: \Latin{audiebas, audiebamus, audiebant}. Future: \Latin{audies, audiemus, audient}.}
\Exercise{Parse and translate \Latin{tollis, custodiebat, glorificabimus, regent}.}{\Latin{tollis}: 2 singular present, \English{you take away}; \Latin{custodiebat}: 3 singular imperfect, \English{he/she was guarding}; \Latin{glorificabimus}: 1 plural future, \English{we shall glorify}; \Latin{regent}: 3 plural future, \English{they will rule}.}
\Exercise{Transform \Latin{Credo in unum Deum} to first-person plural and imperfect singular.}{\Latin{Credimus in unum Deum}; \Latin{Credebam in unum Deum}.}
\Exercise{Diagnose \Latin{Agnus Dei, qui tollit peccata mundi} when the Lamb is addressed as \English{you}.}{\Latin{Agnus Dei, qui tollis peccata mundi}. Relative verb must be 2 singular because the antecedent is addressed directly.}
\PassageExercise{Agnus Dei, qui tollis peccata mundi, miserere nobis.}{Parse the active finite verb fully; identify the command that belongs to a later verb class; translate.}{\Latin{tollis}: 2 singular present active indicative of \Latin{tollo}; subject \Latin{qui}, antecedent \Latin{Agnus}; object \Latin{peccata}. \Latin{miserere}: 2 singular present deponent imperative with active meaning. \English{Lamb of God, who take away the sins of the world, have mercy on us.}}
\end{worksheet}

\end{courseunit}

\begin{courseunit}{I.9}
\chapter{Unit 9: The Perfect Active System}

\begin{worksheet}{I.33}{A}{Unit 9: The Perfect Active System}{Variant A}{Perfect stems, endings, and completed action}{Show the third principal part used for every irregular stem.}
\Exercise{Derive perfect stems from \Latin{laudavi, monui, rexi, cepi, audivi, dixi}.}{\Latin{laudav-, monu-, rex-, cep-, audiv-, dix-}.}
\Exercise{Write all six perfect active indicative forms of \Latin{laudo}.}{\Latin{laudavi, laudavisti, laudavit, laudavimus, laudavistis, laudaverunt}.}
\Exercise{Give 1 singular, 2 plural, and 3 plural pluperfect and future perfect of \Latin{audio}.}{Pluperfect: \Latin{audiveram, audiveratis, audiverant}. Future perfect: \Latin{audivero, audiveritis, audiverint}.}
\Exercise{Contrast \Latin{laudabat} with \Latin{laudavit} in form and aspect.}{\Latin{laudabat}: 3 singular imperfect active indicative, \English{was praising/used to praise}; \Latin{laudavit}: 3 singular perfect active indicative, \English{praised/has praised}. The former views action as ongoing or repeated, the latter as complete.}
\Exercise{Correct \Latin{heri Deum laudabit}, intended \English{yesterday he praised God}.}{\Latin{Heri Deum laudavit}. \Latin{Laudabit} is future; completed past requires perfect \Latin{laudavit}.}
\PassageExercise{Resurrexit, sicut dixit.}{This is an excerpt of the traditional \Latin{Regina cæli} antiphon whose controlling edition remains unresolved; parse both verbs and translate.}{\Latin{resurrexit}: 3 singular perfect active indicative of \Latin{resurgo, resurgere, resurrexi}; \Latin{dixit}: 3 singular perfect active indicative of \Latin{dico, dicere, dixi}. \English{He has risen, as he said.}}
\end{worksheet}

\begin{worksheet}{I.34}{B}{Unit 9: The Perfect Active System}{Variant B}{Perfect-system tense contrasts in sacred text}{Use timeline language in each answer.}
\Exercise{Write the perfect endings and attach them to perfect stem \Latin{fec-}.}{Endings \Latin{-i, -isti, -it, -imus, -istis, -erunt}; forms \Latin{feci, fecisti, fecit, fecimus, fecistis, fecerunt}.}
\Exercise{Parse \Latin{videramus, creaverint, dixistis, resurrexit}.}{\Latin{videramus}: 1 plural pluperfect active indicative; \Latin{creaverint}: 3 plural future perfect active indicative (or perfect subjunctive outside this frame); \Latin{dixistis}: 2 plural perfect active indicative; \Latin{resurrexit}: 3 singular perfect active indicative.}
\Exercise{Transform \Latin{Deus mundum creavit} into pluperfect and future perfect without changing person or number.}{\Latin{Deus mundum creaverat}; \Latin{Deus mundum creaverit}.}
\Exercise{Translate \English{we saw the light}, \English{they had heard the word}, and \English{you (plural) will have praised God}.}{\Latin{Lucem vidimus}; \Latin{verbum audiverant}; \Latin{Deum laudaveritis}.}
\Exercise{Diagnose \Latin{postquam audiverunt, venerant} if the intended sequence is \English{after they had heard, they came}.}{\Latin{Postquam audiverant, venerunt}. Prior completed event uses pluperfect \Latin{audiverant}; subsequent completed event uses perfect \Latin{venerunt}.}
\PassageExercise{Et resurrexit tertia die, secundum Scripturas.}{Parse the finite verb, identify its understood subject and the two adverbial phrases, then translate.}{\Latin{resurrexit}: 3 singular perfect active indicative; understood subject is Christ. \Latin{tertia die}: ablative of time when; \Latin{secundum Scripturas}: accusative after \Latin{secundum}. \English{And he rose again on the third day, according to the Scriptures.}}
\end{worksheet}

\begin{worksheet}{I.35}{C}{Unit 9: The Perfect Active System}{Variant C}{Irregular perfect stems and clause-level parsing}{Do not manufacture a perfect stem from the present stem.}
\Exercise{Match present to perfect and derive the stem: \Latin{facio--feci; dico--dixi; video--vidi; capio--cepi; venio--veni; pono--posui}.}{Stems: \Latin{fec-, dix-, vid-, cep-, ven-, posu-}.}
\Exercise{Give 3 singular perfect, 1 plural pluperfect, and 3 plural future perfect of \Latin{dico}.}{\Latin{dixit, dixeramus, dixerint}.}
\Exercise{Parse and translate \Latin{fecit, dispersit, deposuit, exaltavit}.}{Each is 3 singular perfect active indicative: \English{he did/made; he scattered; he put down; he exalted}.}
\Exercise{Change \Latin{vidi aquam}, \English{I saw water}, to \English{we had seen water} and \English{they will have seen water}.}{\Latin{aquam videramus}; \Latin{aquam viderint}.}
\Exercise{Correct \Latin{Deus fecerunt potentiam}, preserving singular subject.}{\Latin{Deus fecit potentiam}. Singular \Latin{Deus} requires 3 singular \Latin{fecit}, not plural \Latin{fecerunt}.}
\PassageExercise{Fecit potentiam in brachio suo; dispersit superbos mente cordis sui.}{Parse finite verbs, objects, reflexive possessives, and ablative phrases; translate.}{\Latin{fecit, dispersit}: 3 singular perfect active; understood subject God. Objects: \Latin{potentiam} and substantival \Latin{superbos}. In \Latin{brachio suo}, \Latin{suo} refers to God; in the compact phrase \Latin{mente cordis sui}, \Latin{sui} is an indirect reflexive referring to the proud and their own heart. \Latin{in brachio suo} and \Latin{mente} are ablative expressions; \Latin{cordis} is genitive. \English{He has shown might with his arm; he has scattered the proud in the thought of their heart.}}
\end{worksheet}

\begin{worksheet}{I.36}{D}{Unit 9: The Perfect Active System}{Variant D}{Perfect versus imperfect and coordinated completed acts}{Use as a fresh cumulative return after the other forms.}
\Exercise{Write all persons of the pluperfect and future perfect of \Latin{rego} using perfect stem \Latin{rex-}.}{Pluperfect: \Latin{rexeram, rexeras, rexerat, rexeramus, rexeratis, rexerant}. Future perfect: \Latin{rexero, rexeris, rexerit, rexerimus, rexeritis, rexerint}.}
\Exercise{Classify and translate \Latin{creabat, creavit, creaverat, creaverit}.}{3 singular imperfect \English{was creating}; 3 singular perfect \English{created/has created}; 3 singular pluperfect \English{had created}; 3 singular future perfect \English{will have created}.}
\Exercise{Give the alternative third-person plural perfect ending and form it from \Latin{laudav-}; state the preferred production form.}{Alternative ending \Latin{-ere}: \Latin{laudavere}. Preferred production form: \Latin{laudaverunt}.}
\Exercise{Translate \English{he said}, \English{we had done}, \English{they have come}, and \English{I shall have heard}.}{\Latin{dixit; feceramus; venerunt; audivero}.}
\Exercise{Diagnose \Latin{heri videbamus aquam et vidimus templum} if both acts are presented as single completed events.}{\Latin{Heri vidimus aquam et vidimus templum}. Imperfect \Latin{videbamus} would present the first seeing as ongoing/repeated; perfect presents both as completed.}
\PassageExercise{Deposuit potentes de sede, et exaltavit humiles.}{Parse verbs and substantival adjectives; identify the governed phrase and translate.}{\Latin{deposuit, exaltavit}: 3 singular perfect active indicative with understood divine subject. \Latin{potentes, humiles}: masculine/feminine plural accusative substantival adjectives, direct objects. \Latin{de sede}: ablative after \Latin{de}. \English{He has put down the mighty from their seat and exalted the lowly.}}
\end{worksheet}

\end{courseunit}

\begin{courseunit}{I.10}
\chapter{Unit 10: Passive Voice and Deponents}

\begin{worksheet}{I.37}{A}{Unit 10: Passive Voice and Deponents}{Variant A}{Present-system passives, agency, and means}{Keep grammatical subject distinct from personal agent.}
\Exercise{Write all six present passive forms of \Latin{laudo} and \Latin{rego}.}{\Latin{laudor, laudaris, laudatur, laudamur, laudamini, laudantur}; \Latin{regor, regeris, regitur, regimur, regimini, reguntur}.}
\Exercise{Convert without changing person, number, tense, or mood: \Latin{laudat, regimus, audiebant, capietis}.}{\Latin{laudatur, regimur, audiebantur, capiemini}.}
\Exercise{Give the present passive infinitives of \Latin{laudo, moneo, rego, capio}, and \Latin{audio}; identify the special third-conjugation ending.}{\Latin{laudari, moneri, regi, capi, audiri}. Third and third-\Latin{io} use \Latin{-i}, not \Latin{-eri}.}
\Exercise{Translate \Latin{a Domino laudatur} and \Latin{verbo sanatur}; distinguish agent from means.}{\English{He/she is praised by the Lord}: personal agent takes \Latin{a} plus ablative. \English{He/she is healed by a word}: impersonal means uses ablative alone.}
\Exercise{Correct \Latin{animæ custoditur a Domino} when plural souls are the subject.}{\Latin{Animæ custodiuntur a Domino}. Plural nominative subject requires 3 plural passive \Latin{custodiuntur}.}
\PassageExercise{Petite, et dabitur vobis.}{Parse both verbs, identify the subject of the passive and the recipient, then translate.}{\Latin{petite}: 2 plural present active imperative; \Latin{dabitur}: 3 singular future passive indicative, with understood neuter subject \English{it}; \Latin{vobis}: dative plural recipient. \English{Ask, and it will be given to you.}}
\end{worksheet}

\begin{worksheet}{I.38}{B}{Unit 10: Passive Voice and Deponents}{Variant B}{Perfect-system passives and participle agreement}{Make the participle agree with the passive subject, not the agent.}
\Exercise{Construct the perfect passive of \Latin{laudo} for (a) one man, (b) one woman, (c) several men, (d) several women, and (e) several neuter things.}{\Latin{laudatus est; laudata est; laudati sunt; laudatæ sunt; laudata sunt}.}
\Exercise{Parse and translate \Latin{custodita erat, auditi erunt, laudatum est, captæ sunt}.}{\Latin{custodita erat}: feminine singular pluperfect passive, \English{she/it had been guarded}; \Latin{auditi erunt}: masculine plural future perfect passive, \English{they will have been heard}; \Latin{laudatum est}: neuter singular perfect passive, \English{it was praised}; \Latin{captæ sunt}: feminine plural perfect passive, \English{they were captured}.}
\Exercise{Change active \Latin{Dominus animas custodivit} into passive without omitting the agent.}{\Latin{Animæ a Domino custoditæ sunt}. Object becomes nominative subject; agent is \Latin{a Domino}; participle agrees feminine plural.}
\Exercise{Transform \Latin{anima sanata est} to pluperfect and future perfect, keeping the subject singular.}{\Latin{anima sanata erat}; \Latin{anima sanata erit}.}
\Exercise{Diagnose \Latin{verba a Domino dictus sunt}.}{\Latin{Verba a Domino dicta sunt}. Neuter plural subject \Latin{verba} requires neuter plural participle \Latin{dicta}; auxiliary \Latin{sunt} is correctly plural.}
\PassageExercise{Per quem omnia facta sunt.}{Parse the relative phrase, subject, and finite group; distinguish \Latin{fio} from a regular passive of \Latin{facio}; translate.}{\Latin{per quem}: masculine singular accusative after \Latin{per}; \Latin{omnia}: neuter plural nominative subject; \Latin{facta sunt}: neuter plural perfect of \Latin{fio, fieri, factus sum}, \English{came to be/were made}. \English{Through whom all things were made.}}
\end{worksheet}

\begin{worksheet}{I.39}{C}{Unit 10: Passive Voice and Deponents}{Variant C}{Deponent morphology with active meaning}{State both surface morphology and actual meaning.}
\Exercise{Give the principal parts and meaning of \Latin{confiteor}, \Latin{sequor}, and \Latin{misereor}.}{\Latin{confiteor, confiteri, confessus sum}, confess; \Latin{sequor, sequi, secutus sum}, follow; \Latin{misereor, misereri, misertus sum}, have mercy.}
\Exercise{Parse and translate \Latin{confiteor, sequitur, sequebantur, miseremini}; preserve any genuine ambiguity.}{\Latin{confiteor}: 1 singular present deponent indicative, \English{I confess}; \Latin{sequitur}: 3 singular present deponent indicative, \English{he/she follows}; \Latin{sequebantur}: 3 plural imperfect deponent indicative, \English{they were following}; isolated \Latin{miseremini} can be 2 plural present deponent indicative, \English{you have mercy}, or present deponent imperative, \English{have mercy}. Context decides.}
\Exercise{Construct deponent perfects: \English{Mary followed Christ}; \English{the disciples confessed}; \English{the soul had followed the Lord}.}{\Latin{Maria Christum secuta est}; \Latin{discipuli confessi sunt}; \Latin{anima Dominum secuta erat}. Perfect participles agree with deponent subjects, but meanings remain active.}
\Exercise{Explain why \Latin{sequitur} cannot be translated \English{is followed} merely from its ending.}{Its dictionary entry marks \Latin{sequor} as deponent. Deponents use passive-form endings but active meanings, so \Latin{sequitur} means \English{follows}. Lexical classification is essential.}
\Exercise{Correct \Latin{Discipulus a Christo sequitur}, intended \English{The disciple follows Christ}.}{\Latin{Discipulus Christum sequitur}. Deponent \Latin{sequitur} has active meaning and governs accusative object \Latin{Christum}; \Latin{a Christo} would falsely treat Christ as a passive agent.}
\PassageExercise{Confiteor Deo omnipotenti, beatæ Mariæ semper Virgini.}{Parse the deponent, both dative phrases, and their adjectives; translate.}{\Latin{confiteor}: 1 singular present deponent indicative, active meaning. \Latin{Deo omnipotenti}: masculine singular dative; \Latin{beatæ Mariæ semper Virgini}: feminine singular dative with agreeing \Latin{beatæ}, and \Latin{semper} adverb. \English{I confess to almighty God, to blessed Mary ever Virgin.}}
\end{worksheet}

\begin{worksheet}{I.40}{D}{Unit 10: Passive Voice and Deponents}{Variant D}{True passives, deponents, participles, and \Latin{fio}}{Classify every passive-looking form before translating.}
\Exercise{Classify and parse \Latin{laudatur, sequitur, facta sunt, benedictus est}, giving the contextual questions that remain.}{\Latin{laudatur}: 3 singular present true passive, \English{is praised}. \Latin{sequitur}: 3 singular present deponent, \English{follows}. \Latin{facta sunt}: neuter plural perfect of \Latin{fio}, \English{were made/came to be}. \Latin{benedictus est}: masculine singular perfect passive, \English{was/has been blessed}, or adjective plus \Latin{est}, \English{is blessed}, as context decides.}
\Exercise{Give 3 singular present, imperfect, future, perfect, and pluperfect passive of \Latin{audio}.}{\Latin{auditur, audiebatur, audietur, auditus/audita/auditum est, auditus/audita/auditum erat}; participle gender follows the subject.}
\Exercise{Convert \Latin{Deus laudat sanctos} to present passive and \Latin{Deus creavit mundum} to perfect passive.}{\Latin{Sancti a Deo laudantur}; \Latin{Mundus a Deo creatus est}.}
\Exercise{Translate \English{we follow the Lord}, \English{the words will be heard}, and \English{the souls had been guarded}.}{\Latin{Dominum sequimur}; \Latin{verba audientur}; \Latin{animæ custoditæ erant}.}
\Exercise{Diagnose \Latin{Omnia facti sunt per ipsum}.}{\Latin{Omnia facta sunt per ipsum}. Neuter plural subject \Latin{omnia} requires neuter plural \Latin{facta}.}
\PassageExercise{Ite, missa est.}{Parse the imperative and nominal assertion,
state the conventional force of the fixed formula, and translate.}{\Latin{ite}:
irregular 2 plural present active imperative of \Latin{eo}, \English{go}. In
this Late Latin formula, \Latin{missa} is a feminine singular noun meaning
\English{dismissal}, serving as subject or nominal assertion with \Latin{est}:
\English{Go; it is the dismissal}. It is historically related to \Latin{mitto},
but a participial ellipsis is not the sole contextual parse.}
\end{worksheet}

\end{courseunit}

\begin{courseunit}{I.11}
\chapter{Unit 11: Pronouns and Prepositions}

\begin{worksheet}{I.41}{A}{Unit 11: Pronouns and Prepositions}{Variant A}{Personal/reflexive pronouns and preposition government}{For each pronoun, give case, number, and antecedent where applicable.}
\Exercise{Write nominative, genitive, dative, accusative, and ablative of first- and second-person singular pronouns.}{First: \Latin{ego, mei, mihi, me, me}. Second: \Latin{tu, tui, tibi, te, te}.}
\Exercise{Write the four personal pronouns with \Latin{cum}, singular then plural.}{\Latin{mecum, tecum, nobiscum, vobiscum}. The preposition follows and attaches to these pronouns.}
\Exercise{Distinguish \Latin{Petrus matrem suam amat} from \Latin{Petrus matrem eius amat}.}{First: \English{Peter loves his own mother}; reflexive possessive \Latin{suam} refers to subject Peter. Second: \English{Peter loves his/her [someone else's] mother}; \Latin{eius} refers to a non-subject possessor.}
\Exercise{Choose the correct case and form: \Latin{pro (nos)}, \Latin{sine (tu)}, \Latin{ad (ego)}, \Latin{cum (vos)}.}{\Latin{pro nobis}; \Latin{sine te}; \Latin{ad me}; \Latin{vobiscum}.}
\Exercise{Correct \Latin{Dominus cum vobis} to the conventional pronoun-preposition order.}{\Latin{Dominus vobiscum}. With personal pronouns, \Latin{cum} is postpositive and attached.}
\PassageExercise{Dominus vobiscum. Et cum spiritu tuo.}{Parse both prepositional expressions, state the conventional optative construal, and translate.}{\Latin{vobiscum}: second-person plural ablative plus postpositive \Latin{cum}; \Latin{cum spiritu tuo}: ablative noun phrase. An analytic paraphrase may add interpretive \Latin{[sit]} in the first formula, but no verb is printed. \English{The Lord be with you. And with your spirit.}}
\end{worksheet}

\begin{worksheet}{I.42}{B}{Unit 11: Pronouns and Prepositions}{Variant B}{Demonstratives, possessives, and two-case prepositions}{Change every agreeing word when transforming a pronoun.}
\Exercise{Decline neuter singular and plural \Latin{id} from \Latin{is, ea, id}.}{Singular nom./acc. \Latin{id}, gen. \Latin{eius}, dat. \Latin{ei}, abl. \Latin{eo}; plural nom./acc. \Latin{ea}, gen. \Latin{eorum}, dat./abl. \Latin{eis} or \Latin{iis}.}
\Exercise{Give feminine singular and neuter plural of \Latin{hic, hæc, hoc} in nominative, genitive, dative, accusative, and ablative.}{Feminine singular: \Latin{hæc, huius, huic, hanc, hac}. Neuter plural: \Latin{hæc, horum, his, hæc, his}.}
\Exercise{Transform \Latin{per eum, cum eo, in eo} to feminine singular and neuter plural.}{Feminine: \Latin{per eam, cum ea, in ea}. Neuter plural: \Latin{per ea, cum eis, in eis}.}
\Exercise{Translate both contrasts: \English{into this church/in this church}; \English{under that altar/to under that altar}.}{\Latin{in hanc ecclesiam/in hac ecclesia}; \Latin{sub eo altari/sub id altare}. Accusative marks motion toward; ablative marks location.}
\Exercise{Diagnose \Latin{Hoc est Corpus mea}.}{\Latin{Hoc est Corpus meum}. Neuter singular possessive \Latin{meum} agrees with predicate noun \Latin{Corpus}; \Latin{hoc} is neuter singular nominative subject.}
\PassageExercise{Hoc est enim Corpus meum.}{Parse demonstrative, linking verb, predicate noun, and possessive; translate.}{\Latin{hoc}: neuter singular nominative subject; \Latin{est}: 3 singular present; \Latin{Corpus}: neuter singular nominative predicate; \Latin{meum}: agreeing possessive. \English{For this is my Body.}}
\end{worksheet}

\begin{worksheet}{I.43}{C}{Unit 11: Pronouns and Prepositions}{Variant C}{Relative pronouns and antecedent logic}{Take gender and number from the antecedent, case from the relative clause.}
\Exercise{Write the complete singular paradigm of \Latin{qui, quæ, quod}.}{Nom. \Latin{qui, quæ, quod}; gen. \Latin{cuius} all genders; dat. \Latin{cui} all; acc. \Latin{quem, quam, quod}; abl. \Latin{quo, qua, quo}.}
\Exercise{Write the complete plural paradigm of the relative pronoun.}{Nom. \Latin{qui, quæ, quæ}; gen. \Latin{quorum, quarum, quorum}; dat. \Latin{quibus} all; acc. \Latin{quos, quas, quæ}; abl. \Latin{quibus} all.}
\Exercise{Supply the relative:
\begin{enumerate}[label=(\alph*),leftmargin=2em]
  \item \Latin{vir \Blank venit}
  \item \Latin{anima \Blank Deus custodit}
  \item \Latin{verba \Blank audivimus}
  \item \Latin{Patri \Blank gratias agimus}
\end{enumerate}}{\Latin{vir qui venit}; \Latin{anima quam Deus custodit}; \Latin{verba quæ audivimus}; \Latin{Patri cui gratias agimus}.}
\Exercise{For \Latin{dies quam fecit Dominus}, state separately what determines \Latin{quam}'s gender/number and what determines its case.}{Its antecedent \Latin{dies}, feminine in this context and singular, determines feminine singular. Its job as direct object of \Latin{fecit} determines accusative.}
\Exercise{Correct \Latin{Agnus Dei, quem tollis peccata mundi}.}{\Latin{Agnus Dei, qui tollis peccata mundi}. Relative pronoun is subject of \Latin{tollis}, hence nominative \Latin{qui}; \Latin{peccata} is the object.}
\PassageExercise{Pater noster, qui es in cælis, dimitte nobis debita nostra.}{This controlled excerpt joins the opening address directly to a later petition and omits intervening clauses; resolve the relative and parse every phrase.}{\Latin{Pater noster}: masculine singular address; \Latin{qui}: masculine singular nominative subject of \Latin{es}; \Latin{in cælis}: ablative plural; \Latin{nobis}: dative plural recipient; \Latin{debita nostra}: neuter plural accusative object. \English{Our Father, who are in heaven, forgive us our debts.}}
\end{worksheet}

\begin{worksheet}{I.44}{D}{Unit 11: Pronouns and Prepositions}{Variant D}{Core prepositions and cumulative pronoun parsing}{State the governing word for every oblique case.}
\Exercise{Sort by government: \Latin{ad, a/ab, apud, contra, cum, de, e/ex, inter, ob, per, post, pro, propter, secundum, sine, trans}.}{Accusative: \Latin{ad, apud, contra, inter, ob, per, post, propter, secundum, trans}. Ablative: \Latin{a/ab, cum, de, e/ex, pro, sine}.}
\Exercise{Give the correct phrase: \English{through us}; \English{for her}; \English{without them} masculine; \English{from this altar}; \English{toward those churches}.}{\Latin{per nos}; \Latin{pro ea}; \Latin{sine eis}; \Latin{ab hoc altari}; \Latin{ad eas ecclesias}.}
\Exercise{Parse the repeated pronoun in \Latin{per ipsum, cum ipso, in ipso}; explain each case.}{\Latin{ipsum}: masculine singular accusative after \Latin{per}; both \Latin{ipso} forms: masculine singular ablative after \Latin{cum} and after \Latin{in} denoting state or location. Antecedent is Christ.}
\Exercise{Explain why \Latin{nostram} in \Latin{propter nostram salutem} is not a form of personal pronoun \Latin{nos}.}{It is the feminine singular accusative possessive adjective \Latin{noster, nostra, nostrum}, agreeing with \Latin{salutem}; \Latin{nos} itself is the indeclinable-looking accusative personal form in \Latin{propter nos}.}
\Exercise{Correct \Latin{de cælos, propter nobis, cum eos}.}{\Latin{de cælis}, ablative after \Latin{de}; \Latin{propter nos}, accusative after \Latin{propter}; \Latin{cum eis}, ablative after \Latin{cum}.}
\PassageExercise{Propter nos homines et propter nostram salutem descendit de cælis.}{Parse both accusative phrases and the source phrase; identify the finite verb and translate.}{\Latin{propter nos homines}: accusative plural, with \Latin{homines} in apposition to \Latin{nos}; \Latin{propter nostram salutem}: feminine singular accusative; \Latin{de cælis}: masculine plural ablative; \Latin{descendit}: 3 singular perfect active in the Creed's narrative. \English{For us men and for our salvation he came down from heaven.}}
\end{worksheet}

\end{courseunit}

\begin{courseunit}{I.12}
\chapter{Unit 12: Consolidated Missal Reading}

\begin{worksheet}{I.45}{A}{Unit 12: Consolidated Missal Reading}{Variant A}{A complete seven-pass analysis of the Ordinary}{Work without a prepared translation; cite morphology for each conclusion.}
\Exercise{Write the seven reading passes in operational form, one action per line.}{(1) Mark punctuation/conjunctions and bracket clauses. (2) Underline and parse finite verbs; bracket only supplied words. (3) Find stated or understood subjects. (4) Give dictionary form, gender, case, number, and function for nouns. (5) Draw agreement/antecedent links. (6) Box prepositional phrases and identify genitive/dative relations. (7) Write close, then natural, English and record departures.}
\PassageExercise{Gloria in excelsis Deo, et in terra pax hominibus bonæ voluntatis.}{State the conventional construal, parse every nominal phrase, and give close English.}{This is a coordinated verbless acclamation conventionally read as a wish; an analytic paraphrase may add interpretive \Latin{[sit]} with \Latin{gloria} and \Latin{pax}. Both are feminine singular nominative; \Latin{in excelsis} is neuter plural ablative substantival; \Latin{Deo} dative recipient; \Latin{in terra} ablative; \Latin{hominibus} dative plural; \Latin{bonæ voluntatis} feminine singular genitive. \English{Glory [be] to God in the highest, and on earth peace [be] to people of good will.}}
\PassageExercise{Sanctus, Sanctus, Sanctus Dominus Deus Sabaoth. Pleni sunt cæli et terra gloria tua.}{Account for agreement, apposition, indeclinability, subject-verb number, and the ablative; translate.}{\Latin{Sanctus} is masculine singular nominative in repeated acclamatory/predicate relation to \Latin{Dominus}; \Latin{Deus} is appositional; \Latin{Sabaoth} indeclinable. \Latin{cæli et terra} is compound subject of plural \Latin{sunt}; predicate \Latin{pleni} is masculine plural; \Latin{gloria tua} ablative after \Latin{pleni}. \English{Holy, Holy, Holy Lord God of hosts. Heaven and earth are full of your glory.}}
% This cumulative form deliberately occupies two balanced pages in both the
% learner packet and key; do not leave only its final passage on a tail page.
\clearpage
\Exercise{Transform active \Latin{Deus animas custodivit} to passive, then make the passive subject singular; translate both results.}{Plural: \Latin{Animæ a Deo custoditæ sunt}, \English{The souls were guarded by God}. Singular: \Latin{Anima a Deo custodita est}, \English{The soul was guarded by God}.}
\Exercise{Find and correct every error: \Latin{Domine, non sumus dignum; sed dic verba, et animæ nostra sanabitur}. Assume one speaker and one soul.}{\Latin{Domine, non sum dignus; sed dic verbo, et anima nostra sanabitur}. First-person singular \Latin{sum}; nominative predicate \Latin{dignus}; ablative \Latin{verbo}; feminine singular nominative \Latin{anima nostra} with singular passive \Latin{sanabitur}.}
\PassageExercise{Domine, non sum dignus ut intres sub tectum meum: sed tantum dic verbo, et sanabitur anima mea.}{Parse every finite form and explain the \Latin{ut}-clause and the cases in the \Latin{sub} phrase, with \Latin{dic}, and with the passive; translate.}{\Latin{sum}: 1 singular present active indicative; \Latin{intres}: 2 singular present active subjunctive in an \Latin{ut}-clause completing \Latin{dignus}; \Latin{dic}: 2 singular present active imperative; \Latin{sanabitur}: 3 singular future passive indicative. \Latin{sub tectum meum}: accusative motion; \Latin{verbo}: ablative means; \Latin{anima mea}: nominative subject. \English{Lord, I am not worthy that you should enter under my roof; but only say the word, and my soul shall be healed.}}
\end{worksheet}

\begin{worksheet}{I.46}{B}{Unit 12: Consolidated Missal Reading}{Variant B}{Cumulative reading with pronouns and voice contrasts}{Give close English before idiomatic English.}
\PassageExercise{In principio erat Verbum, et Verbum erat apud Deum, et Deus erat Verbum.}{Parse each \Latin{erat}; distinguish subjects and predicates without relying on position; translate.}{Every \Latin{erat} is 3 singular imperfect active indicative of \Latin{sum}. First subject \Latin{Verbum}; second subject \Latin{Verbum} with \Latin{apud Deum} accusative prepositional phrase. In the third, \Latin{Verbum} is contextual subject and \Latin{Deus} predicate nominative. \English{In the beginning was the Word, and the Word was with God, and the Word was God.}}
\PassageExercise{Hoc erat in principio apud Deum. Omnia per ipsum facta sunt.}{Resolve \Latin{hoc/ipsum}, parse both finite groups, and translate.}{\Latin{hoc}: neuter singular nominative demonstrative referring to the Word; \Latin{erat}: 3 singular imperfect. \Latin{apud Deum}: accusative. \Latin{omnia}: neuter plural nominative; \Latin{per ipsum}: masculine singular accusative referring to the Word; \Latin{facta sunt}: neuter plural perfect of \Latin{fio}. \English{He was in the beginning with God. All things were made through him.}}
\Exercise{Compose and fully label \English{The holy souls had heard the eternal Word}.}{\Latin{Animæ sanctæ Verbum æternum audiverant}. \Latin{animæ sanctæ}: feminine plural nominative subject; \Latin{Verbum æternum}: neuter singular accusative object; \Latin{audiverant}: 3 plural pluperfect active indicative.}
\LearnerWorksheetBreak
\Exercise{Correct \Latin{Per eo omnia factus est}; then explain every changed coordinate.}{\Latin{Per eum omnia facta sunt}. \Latin{per} takes masculine singular accusative \Latin{eum}; neuter plural subject \Latin{omnia} requires neuter plural \Latin{facta} and plural \Latin{sunt}.}
\PassageExercise{Agnus Dei, qui tollis peccata mundi, miserere nobis.}{Resolve the nominative used in address, relative syntax, genitives, object, and deponent; translate.}{\Latin{Agnus} is nominative singular used for direct address; \Latin{Dei} genitive. \Latin{qui} is nominative subject of 2 singular present active \Latin{tollis}; \Latin{peccata} accusative plural with genitive \Latin{mundi}. \Latin{miserere}: 2 singular present deponent imperative, active meaning; \Latin{nobis} its liturgical complement. \English{Lamb of God, who take away the sins of the world, have mercy on us.}}
\PassageExercise{Agnus Dei, qui tollis peccata mundi, dona nobis pacem.}{Compare the command's syntax with the preceding deponent petition and translate.}{The relative clause is unchanged. Active imperative \Latin{dona} is 2 singular and takes dative recipient \Latin{nobis} plus accusative object \Latin{pacem}; unlike deponent \Latin{miserere}, it is active in both form and meaning. \English{Lamb of God, who take away the sins of the world, grant us peace.}}
\end{worksheet}

\begin{worksheet}{I.47}{C}{Unit 12: Consolidated Missal Reading}{Variant C}{Psalm and Creed synthesis}{Use dictionary entries only for meanings and principal parts.}
\PassageExercise{Emitte lucem tuam et veritatem tuam.}{Parse the command and both coordinated objects with agreement; translate.}{\Latin{emitte}: 2 singular present active imperative of third-conjugation \Latin{emitto}; \Latin{lucem tuam} and \Latin{veritatem tuam}: feminine singular accusative objects with agreeing possessives. \English{Send forth your light and your truth.}}
\PassageExercise{Ipsa me deduxerunt et adduxerunt in montem sanctum tuum, et in tabernacula tua.}{Resolve \Latin{ipsa}, parse verbs, object, both directional phrases, and translate.}{\Latin{ipsa}: neuter plural nominative referring collectively to light and truth; \Latin{deduxerunt/adduxerunt}: 3 plural perfect active; \Latin{me}: accusative object. Both \Latin{in} phrases are directional accusatives: \Latin{montem sanctum tuum} masculine singular and \Latin{tabernacula tua} neuter plural. \English{They have led me and brought me unto your holy mountain and into your tabernacles.}}
\LearnerWorksheetBreak
\Exercise{For \Latin{creavit, descendit, resurrexit, facta sunt}, give principal part evidence, full parse, and a close rendering.}{\Latin{creavit}: from \Latin{creavi}, 3 singular perfect active, \English{created}; \Latin{descendit}: from \Latin{descendi}, 3 singular perfect active in Creed context, \English{came down}; \Latin{resurrexit}: from \Latin{resurrexi}, 3 singular perfect active, \English{rose}; \Latin{facta sunt}: from \Latin{fio, factus sum}, neuter plural perfect, \English{were made}.}
\Exercise{Compose \English{For our salvation the Lord came from heaven}; then mark preposition government and verb tense.}{\Latin{Propter salutem nostram Dominus de cælo venit}. \Latin{propter} governs feminine singular accusative \Latin{salutem nostram}; \Latin{de} governs neuter singular ablative \Latin{cælo}; \Latin{venit} is 3 singular perfect active in this completed-event context.}
\Exercise{Correct \Latin{In nominem Patrem et Filius et Spiritum Sanctum}.}{\Latin{In nomine Patris, et Filii, et Spiritus Sancti}. \Latin{in} takes ablative \Latin{nomine}; each named possessor is genitive, with \Latin{Sancti} agreeing with \Latin{Spiritus}.}
\PassageExercise{Et resurrexit tertia die, secundum Scripturas: et ascendit in cælum.}{Parse completed verbs and all adverbial case phrases; translate.}{\Latin{resurrexit, ascendit}: 3 singular perfect active indicatives with understood Christ as subject. \Latin{tertia die}: ablative time when; \Latin{secundum Scripturas}: accusative after \Latin{secundum}; \Latin{in cælum}: accusative motion toward. \English{And he rose again on the third day, according to the Scriptures, and ascended into heaven.}}
\end{worksheet}

\begin{worksheet}{I.48}{D}{Unit 12: Consolidated Missal Reading}{Variant D}{Independent cumulative reading}{Work through the passage without the model, then explain every correction by rule.}
\PassageExercise{Laudamus te, benedicimus te, adoramus te, glorificamus te, gratias agimus tibi propter magnam gloriam tuam.}{Parse all verbs, pronouns, objects, recipient, and prepositional phrase; translate.}{The first four verbs are 1 plural present active with accusative \Latin{te}. \Latin{agimus} is 1 plural present active; \Latin{gratias} its accusative plural object and \Latin{tibi} dative recipient. \Latin{propter} governs feminine singular accusative \Latin{magnam gloriam tuam}. \English{We praise you, bless you, adore you, glorify you; we give you thanks because of your great glory.}}
\PassageExercise{Benedictus qui venit in nomine Domini.}{Explain the verbless acclamation, resolve the relative, parse the phrase, and translate.}{\Latin{Benedictus} is a substantival masculine singular nominative, and \Latin{qui venit} identifies the blessed one; conventional English supplies a copula, but none is printed. \Latin{Qui} is masculine singular nominative subject of 3 singular present \Latin{venit}; \Latin{in nomine} is ablative; \Latin{Domini} genitive. \English{Blessed is he who comes in the name of the Lord.}}
\Exercise{Transform \Latin{Verbum a Domino auditur} into imperfect passive, future passive, perfect passive, and active present.}{\Latin{Verbum a Domino audiebatur}; \Latin{Verbum a Domino audietur}; \Latin{Verbum a Domino auditum est}; active \Latin{Dominus Verbum audit}.}
\LearnerWorksheetBreak
\Exercise{Compose \English{All holy names will have been praised by the faithful}; parse your finite group.}{\Latin{Omnia nomina sancta a fidelibus laudata erunt}. \Latin{laudata erunt}: neuter plural future perfect passive indicative agreeing with \Latin{omnia nomina sancta}; \Latin{a fidelibus} personal agent.}
\Exercise{Diagnose and correct all errors: \Latin{Pater nostra, qui est in cælos, dona nos pacis}.}{\Latin{Pater noster, qui es in cælis, dona nobis pacem}. Masculine singular address takes \Latin{noster}; addressed Father requires 2 singular \Latin{es}; \Latin{in} denoting place where takes ablative plural \Latin{cælis}; recipient is dative \Latin{nobis}; object is accusative \Latin{pacem}.}
\PassageExercise{Requiem æternam dona eis, Domine, et lux perpetua luceat eis.}{Provide a complete morphological map, identify the two petitions, and translate.}{First petition: 2 singular active imperative \Latin{dona}, object \Latin{requiem æternam} feminine singular accusative, recipient \Latin{eis}, address \Latin{Domine}. Second: \Latin{lux perpetua} feminine singular nominative subject; \Latin{luceat} 3 singular present active subjunctive of wish; \Latin{eis} dative. \English{Grant them eternal rest, O Lord, and may perpetual light shine upon them.}}
\end{worksheet}
\end{courseunit}
